% !TEX root = ../main.tex 
\clearpage 
\pagestyle{empty} 
\newgeometry{left=0mm,right=0mm,top=0mm,bottom=0mm}
% \begin{tcolorbox}[ 
%     enhanced, 
%     colback=white, 
%     colframe=black, 
%     boxrule=0.6pt, 
%     sharp corners, 
%     width=\textwidth, 
%     height=\textheight, 
%     valign=top, 
%     left=12mm, 
%     right=12mm, 
%     top=10mm, 
%     bottom=10mm 
% ] 

\begin{tcolorbox}[
    enhanced,
    colback=white,
    colframe=black,
    boxrule=0.6pt,
    sharp corners,
    width=\paperwidth,
    height=\paperheight,
    valign=top,
    left=25mm,
    right=25mm,
    top=20mm,
    bottom=20mm
]
\centering 
{\Large\bfseries NHẬN XÉT CỦA GIẢNG VIÊN HƯỚNG DẪN\par} 
\vspace{1cm} 

\begin{flushleft} 
\large 
\textbf{1. Mục tiêu và nội dung của đồ án} 

\vspace{0.4cm} 
(a) Mục tiêu: 
\par\dotfill 
\par\dotfill 
\par\dotfill 

\vspace{0.6cm} 
(b) Nội dung: 
\par\dotfill 
\par\dotfill 
\par\dotfill 
\par\dotfill 

\vspace{0.8cm} 
\textbf{2. Kết quả đạt được:} 
\par\dotfill 
\par\dotfill 
\par\dotfill 
\par\dotfill 

\vspace{0.8cm} 
\textbf{3. Ý thức làm việc của sinh viên:} 
\par\dotfill 
\par\dotfill 
\par\dotfill 
\par\dotfill 

\end{flushleft} 

\vfill 

\begin{flushright} 
\textit{Hà Nội, ngày \hspace{1cm} tháng \hspace{1cm} năm 2026} 

\textbf{Giảng viên hướng dẫn} 
\vspace{1.5cm} 
\end{flushright} 
\end{tcolorbox}
\restoregeometry
